\documentclass[aps,prx,reprint,nofootinbib,floatfix]{revtex4-2}
\usepackage{amsmath,amssymb,amsthm}
\usepackage{graphicx}
\usepackage{hyperref}
\usepackage{xcolor}
\usepackage{listings}
\usepackage{natbib}

%% ═══════════════════════════════════════════════════════════════════
%% \figuredeferred{id}{reason} — the declared-deferral form
%% ═══════════════════════════════════════════════════════════════════
%%
%% WAVE_EXECUTION_PIPELINE.md §Stage 9: "A figure that is planned but not yet
%% drawn is declared with \figuredeferred{id}{reason}, never omitted silently."
%% The law mandated this form and, measured 2026-08-09, it had ZERO uses across
%% all 21 bundle drafts — nine of which carried no figures at all. Nine bundles
%% with no plan is a charter defect; nine bundles with a *stated* plan and a
%% stated blocker is an honest manuscript in progress, and a referee can tell
%% the two apart only if the plan is on the page.
%%
%% The deferral is DELIBERATELY VISIBLE in the compiled PDF. A silent macro
%% would satisfy `bundle_figure_adequacy` while showing a reader nothing, which
%% is the vacuous-Stage-9 failure ("ALL figures PASS" over an empty set) rebuilt
%% one layer up.
%%
%% ⚠️ The reason must name what BLOCKS the figure, not that it is absent.
%% "not yet drawn" is boilerplate and defeats the purpose; "awaiting the L=8
%% m→0 production run" is a blocker a reader can check and a later pass can
%% discharge.
%%
%% Definition lives here, beside docs/counts.tex, because 21 drafts consume it
%% and a per-draft copy is 21 places for the rendering to drift.

%% Rendered as a non-floating box on purpose. A `figure` float would be subject
%% to revtex's two-column placement rules across eleven drafts for no gain, and
%% a deferral does not need a figure number: it is a promise, not a result.

\providecommand{\figuredeferred}[2]{%
  \par\medskip\noindent
  \fbox{\parbox{0.92\columnwidth}{%
    \small\textbf{Figure deferred: \texttt{#1}.}
    Planned for this section; not yet produced. \emph{Blocked on:} #2.%
  }}%
  \par\medskip
}


\newtheorem{thm}{Theorem}
\newtheorem{lem}[thm]{Lemma}
\newtheorem{defn}[thm]{Definition}
\newtheorem{prop}[thm]{Proposition}

\begin{document}

\title{A machine-checked finite-horizon convergence theorem for belief
propagation, and why 4-cycle density does not screen for it}

\author{John G.~Roehm}
\affiliation{Independent research, USA}

\date{\today}

\begin{abstract}
We prove, in the Lean~4 proof assistant and with kernel-only axiom closure,
that synchronous belief propagation on a finite factor graph reaches a genuine
fixed point of the message-update map within a bounded number of sweeps, from
any initial message bundle, whenever the directed message endpoints carry a
topological rank certificate. The proof factors through a reusable statement
about arbitrary rank-local iterations on a finite coordinate type, and is
instantiated on the actual belief-propagation update rather than on a proxy:
rank locality is discharged by reading off the cavity neighbourhoods of the
variable and factor updates. We then examine the combinatorial statistic the
same development builds to detect that regime, a large-deviation rate function
whose Bernoulli parameter is the 4-cycle density of the factor graph, and
report a negative result. The rate admits a closed Legendre evaluation
$-\log(1-p)$, vanishes exactly on four-cycle-free graphs, is strictly positive
and strictly monotone off them, and takes the kernel-checked value
$\log(4/3)$ on the complete bipartite $2{\times}2$ graph; but four-cycle-freeness
is implied by, and does not imply, acyclicity, and it holds automatically for
every model whose factors couple at most two variables on a simple interaction
graph. The rate is therefore identically zero across the lattice families of
the tensor-network simulations that motivate it, including those whose loop
corrections are the technical core of that work, and its normalization by the
ambient tuple space makes a fixed loop motif's density fall as the system
grows. One step separates the convergence theorem from an unconditional
statement about trees: the existence of a rank certificate on an acyclic
bipartite incidence graph is assumed, not derived. We state that obligation
explicitly, along with the three others that a genuine simulability criterion
would require.
\end{abstract}

\maketitle

\section{Introduction}
\label{sec:intro}

Two recent results narrowed the practical gap between classical and quantum
simulation of correlated many-body systems, and they are the reason a formal
criterion for where that gap lies is worth stating precisely.

Tindall, Mello, Fishman, Stoudenmire and
Sels~\cite{TindallMelloFishmanStoudenmireSels2026Science392} simulated the
glassy annealing dynamics of a disordered transverse-field Ising model on
cylindrical, diamond cubic and dimerized cubic lattices using tensor networks
whose connectivity mirrors the lattice, evolved by a belief-propagation-based
simple-update scheme and measured with controlled contraction. On the
cylindrical and diamond lattices their errors fall noticeably below those of
the D-Wave Advantage2 annealer on the same instances; on the dimerized cubic
lattice the two are comparable. Separately, they simulated a system of over
$300$ qubits and used the resulting scaling collapse of the correlation
function to extract a Kibble--Zurek exponent $\mu \approx 2.70$ to $2.75$,
consistent with independent Monte Carlo estimates. Two features of that work
matter for what follows. The first is that plain belief propagation was not
accurate enough: the authors report that standard belief-propagation
contraction, successful on tree-like systems, does not suffice for the lattices
they study, and they use loop-corrected belief propagation and matrix-product-state
message passing instead. The second is that they attribute a specific error
increase in the dimerized cubic geometry to a loop of size three created by the
periodic boundary, and remove inherent loops by pairing dimers into single
tensors so that the structure becomes ``more tree-like''. Loop content, at a
specific size, is the operative variable.

Antão, Sun, Fumega and Lado~\cite{AntaoSunFumegaLado2026PRL136} computed local
topological markers on quasicrystalline systems of up to $268$ million sites.
Their method represents the ground-state projector $\hat{\mathcal P}$ as a
matrix-product operator obtained from a Chebyshev expansion of the Hamiltonian
MPO, then builds the Bianco--Resta local Chern
marker~\cite{BiancoResta2011}
$C_\alpha = 2\pi i \langle \alpha | \hat Q \hat X \hat P \hat Y \hat Q -
\hat P \hat X \hat Q \hat Y \hat P | \alpha\rangle$
as a tensor-network contraction. Applied to a $\pi$-flux square lattice with an
incommensurate $C_8$ or $C_{10}$ modulation of the second-neighbour imaginary
hopping, the marker resolves spatially alternating Chern domains of width on the
order of $40$ atoms.

A demarcation criterion separating these regimes would have to be a statement
about the interaction structure that predicts, in advance, whether the classical
method suffices. This paper does not supply one. What it supplies is a
machine-checked theorem about when the message-passing iteration itself
terminates, and a machine-checked account of why the combinatorial statistic we
built to detect that condition does not do the job it was named for.

\subsection{Contributions}

\begin{itemize}
\item A formalization of belief propagation on bipartite factor graphs in the
Yedidia--Freeman--Weiss convention~\cite{YedidiaFreemanWeiss2003}, including
the message updates, their cavity neighbourhoods, the fixed-point predicate,
beliefs, and a Bethe free energy (Sec.~\ref{sec:bp}).

\item A finite-horizon convergence theorem for synchronous belief propagation
(Sec.~\ref{sec:convergence}). Its mathematical core is a self-contained
statement about any iteration on a finite coordinate type whose one-round
update at a coordinate reads only coordinates of strictly smaller rank: such an
iteration freezes coordinate by coordinate and reaches a genuine fixed point
after $\sup(\mathrm{rank})+1$ rounds. The belief-propagation instantiation
supplies rank locality from the actual update rule, and the rank certificate
that supplies it is proven to force four-cycle-freeness, so it is not a vacuous
hypothesis. An explicit certificate on a three-node tree witnesses that the
theorem fires.

\item A large-deviation rate function on finite factor graphs
(Sec.~\ref{sec:rate}), constructed as the Legendre transform at zero deviation
of the log-moment-generating function of a Bernoulli loop-presence observable
whose success parameter is the graph's 4-cycle density. The closed form
$-\log(1-p)$ is proven by an explicit least-upper-bound argument, and the
resulting object is strictly monotone in loop density rather than
$\{0,1\}$-valued.

\item A negative result about that rate function (Sec.~\ref{sec:limits}). It is
the paper's main finding, and it disqualifies the criterion the development was
built to support.
\end{itemize}

\subsection{What is not claimed}

No theorem here states that a system is classically simulable, or that a
quantum processor holds an advantage on it. No theorem bounds a runtime, an
approximation error, or a bond dimension. The convergence theorem is
conditional on a supplied rank certificate and is silent about which graphs
admit one. The Bethe free energy is defined and given boundary and congruence
properties; it is not shown to be stationary at belief-propagation fixed
points. Nothing in the development computes, approximates or bounds a Chern
number, a local Chern marker, or any other topological invariant.

\section{Belief propagation on factor graphs}
\label{sec:bp}

This section fixes the carriers and the update rule that
Sec.~\ref{sec:convergence} proves a theorem about. It is deliberately thin:
only the objects consumed downstream are developed here, and the remaining
formalized material is summarized at the end of the section and listed in
Appendix~\ref{app:statements}.

\subsection{Carriers}

A factor graph is a variable index type $\nu$, a factor index type $\alpha$,
and a decidable incidence relation.

\begin{defn}[\texttt{FactorGraph}]
\texttt{FactorGraph}\,$\nu\ \alpha$ is a structure with a single field
$\mathrm{incidence} : \alpha \to \nu \to \mathrm{Bool}$. Factor $a$ touches
variable $v$ when $\mathrm{incidence}\ a\ v = \mathtt{true}$.
\end{defn}

The variable state space $X$ is a separate parameter, so the same graph carries
message bundles over different alphabets. A message bundle is a pair of
real-valued families indexed by directed endpoints.

\begin{defn}[\texttt{BPMessages}]
\texttt{BPMessages}\,$\nu\ \alpha\ G\ X$ is a structure with fields
$\mathrm{varToFactor} : \nu \to \alpha \to X \to \mathbb{R}$ and
$\mathrm{factorToVar} : \alpha \to \nu \to X \to \mathbb{R}$.
\end{defn}

Messages are real-valued and unnormalized. Normalization is available as a
predicate (\texttt{IsNormalizedAt}, \texttt{BPNormalized}) rather than built
into the carrier, so that the update rule is a total function and no proof
obligation is incurred at every application.

\subsection{The update rule}

The cavity neighbourhoods are the two \texttt{Finset}s that make the update
local:
\begin{align}
\mathrm{otherFactors}\ G\ v\ a &= \{\, b : b \neq a,\ \mathrm{incidence}\ b\ v \,\},\\
\mathrm{otherVars}\ G\ a\ v &= \{\, u : u \neq v,\ \mathrm{incidence}\ a\ u \,\}.
\end{align}
In the Yedidia--Freeman--Weiss convention~\cite{YedidiaFreemanWeiss2003} the
two unnormalized updates are
\begin{align}
\label{eq:varupdate}
m'_{v \to a}(x) &= \!\!\prod_{b \,\in\, \mathrm{otherFactors}\ G\ v\ a}\!\! m_{b \to v}(x),\\
\label{eq:facupdate}
m'_{a \to v}(x) &= \sum_{\substack{y : \nu \to X \\ y(v) = x}} f_a(y)
  \!\!\prod_{u \,\in\, \mathrm{otherVars}\ G\ a\ v}\!\! m_{u \to a}(y(u)),
\end{align}
formalized as \texttt{bpVariableUpdate} and \texttt{bpFactorUpdate}. The empty
product convention gives the leaf boundary directly: a variable whose only
incident factor is $a$ sends the constant message $1$
(\texttt{bpVariableUpdate\_eq\_one\_of\_no\_other\_factors}). One synchronous
sweep \texttt{bpUpdate} applies both rules at every endpoint simultaneously,
and \texttt{IsBPFixedPoint}\,$m\ f$ is the proposition
$\texttt{bpUpdate}\ m\ f = m$, an equality of message bundles rather than a
numerical tolerance.

\subsection{A modelling caveat in the factor update}
\label{sec:caveat-factor}

Equation~\eqref{eq:facupdate} carries a deviation from the textbook rule that
a reader must be told about. The factor weight is typed
$f : \alpha \to (\nu \to X) \to \mathbb{R}$, a function of a complete
assignment to \emph{all} variables, and the sum accordingly ranges over all
$y : \nu \to X$ with $y(v) = x$ rather than over assignments to $\partial a$.
The product is correctly restricted to the cavity $\partial a \setminus \{v\}$,
so locality of the message in the \emph{messages} is exact, which is what
Sec.~\ref{sec:convergence} needs. Locality in the \emph{weights} is not: if
$f_a$ depends on a variable outside $\partial a$, the message
$m_{a \to v}$ depends on it too. When $f_a$ does factor through
$\partial a$, the sum overcounts the textbook sum by the constant
$|X|^{|\nu| - |\partial a|}$, which rescales each factor message by a
per-endpoint constant. Since the development works with unnormalized messages
and states fixed points as exact equalities, that rescaling is not neutral for
\texttt{IsBPFixedPoint}, and no theorem here relates the fixed points of
\eqref{eq:facupdate} to those of the restricted rule.

\subsection{Elementary properties, and what is absent}

Non-negativity and strict positivity propagate through both updates
(\texttt{bpVariableUpdate\_nonneg\_of\_factor\_msgs\_nonneg},
\texttt{bpVariableUpdate\_pos\_of\_factor\_msgs\_pos},
\texttt{bpFactorUpdate\_nonneg\_of\_inputs\_nonneg}); the factor update is
linear in the factor weight
(\texttt{bpFactorUpdate\_smul\_factor\_weight}) and propagates a hard zero
(\texttt{bpFactorUpdate\_eq\_zero\_of\_factor\_weight\_zero}); the variable
update factorizes out any single cavity message
(\texttt{bpVariableUpdate\_factorization}); and a bundle is a fixed point iff
every iterate reproduces it
(\texttt{IsBPFixedPoint\_iff\_all\_iterates\_invariant}).

The development additionally defines variable and factor beliefs, a Shannon
entropy on the belief simplex, a factor energy with a relative-entropy
decomposition, and a Bethe free energy with a vacuum boundary value and a
congruence property. None of that is consumed by any result in this paper.
It is stated in Appendix~\ref{app:statements} and flagged here for what it
lacks: there is no theorem asserting that the Bethe free energy is stationary
at belief-propagation fixed points, which is the property that makes the Bethe
construction interesting. Presenting the free-energy layer as support for the
convergence theorem would be presenting inventory as argument.

\figuredeferred{fig:d7-incidence-and-cycles}{a factor graph beside its bipartite incidence graph, with a 4-cycle and the 6-cycle that separates four-cycle-freeness from acyclicity; blocked only on drawing time and is the first figure to land}

\section{Rank-local iteration and finite-horizon convergence}
\label{sec:convergence}

Belief propagation on a tree converges in a number of sweeps set by the
diameter, and the reason is a topological ordering: a message is determined by
the messages behind it, and on an acyclic graph ``behind'' is a well-founded
relation. This section proves that statement in the form the argument actually
needs, separating the order-theoretic content from the belief-propagation
content so that the first is reusable and the second is checked against the
real update rule.

\subsection{The abstract engine}

Let $E$ be a finite type, $\beta$ any type, $\mathrm{rank} : E \to \mathbb{N}$,
and $\mathrm{step} : (E \to \beta) \to (E \to \beta)$.

\begin{defn}[\texttt{RankLocalStep}]
$\mathrm{step}$ is \emph{rank-local} for $\mathrm{rank}$ when, for all states
$s, s'$ and every coordinate $e$, agreement of $s$ and $s'$ on every coordinate
of rank strictly below $\mathrm{rank}\ e$ forces
$\mathrm{step}\ s\ e = \mathrm{step}\ s'\ e$.
\end{defn}

The hypothesis is exactly ``one round at $e$ reads only strictly lower ranks'',
stated extensionally so that no syntactic dependency analysis is required.

\begin{thm}[\texttt{rankLocalStep\_coord\_stable}]
\label{thm:coordstable}
If $\mathrm{step}$ is rank-local then for every $s$, every $e$, and every
$n \geq \mathrm{rank}\ e + 1$,
$\mathrm{step}^{[n]}\, s\, e = \mathrm{step}^{[\mathrm{rank}\, e + 1]}\, s\, e$.
\end{thm}

The proof is a strong induction on $\mathrm{rank}\ e$. Writing
$n = m+1$ and pushing one application to the outside on both sides reduces the
goal to an application of rank locality, whose hypothesis is discharged by the
induction hypothesis at each strictly smaller rank: coordinate $e'$ has already
frozen by round $\mathrm{rank}\ e' + 1$, and both $m$ and $\mathrm{rank}\ e$
exceed that bound. Rank-zero coordinates freeze after one round because their
rank-locality hypothesis is vacuous, and every higher coordinate freezes exactly
one round after everything it reads.

\begin{defn}
$\mathrm{convergenceHorizon}\ \mathrm{rank} = \bigl(\sup_{e \in E} \mathrm{rank}\ e\bigr) + 1$.
\end{defn}

\begin{thm}[\texttt{rankLocalStep\_converges}]
\label{thm:converges}
If $\mathrm{step}$ is rank-local then for every initial state $s$, writing
$H = \mathrm{convergenceHorizon}\ \mathrm{rank}$: every iterate past $H$ equals
the $H$-th iterate, and the $H$-th iterate is a fixed point of
$\mathrm{step}$.
\end{thm}

Both conclusions matter and neither implies the other cheaply. The first is
stabilization; the second is that the stable point is genuinely fixed, obtained
by evaluating the stabilization statement one round further and commuting the
extra application to the outside.

\subsection{Instantiation on the real update}

Flattening carries the message bundle onto a single coordinate type. Let
$\mathrm{MsgEndpoint}\ \nu\ \alpha = (\nu \times \alpha) \oplus (\alpha \times \nu)$,
with $\mathrm{inl}(v,a)$ the $v \to a$ endpoint and $\mathrm{inr}(a,v)$ the
$a \to v$ endpoint, and let $\mathrm{flattenMsg}$ and $\mathrm{unflattenMsg}$
be the evident mutually inverse maps to $\mathrm{MsgEndpoint} \to (X \to \mathbb{R})$.
Setting $\mathrm{bpStep}\ G\ f = \mathrm{flattenMsg} \circ (\cdot \mapsto
\texttt{bpUpdate}\ \cdot\ f) \circ \mathrm{unflattenMsg}$ makes one synchronous
sweep an endomorphism of the flattened state, and
\texttt{flattenMsg\_iterate} transports iterates in both directions.

\begin{defn}[\texttt{BPRankCert}]
\label{def:cert}
A rank certificate for $G$ is a rank $r$ on $\mathrm{MsgEndpoint}\ \nu\ \alpha$
with two strict inequalities:
\begin{align}
\label{eq:vargt}
b \neq a,\ \mathrm{incidence}\ b\ v &\ \Rightarrow\ r(\mathrm{inr}(b,v)) < r(\mathrm{inl}(v,a)),\\
\label{eq:facgt}
u \neq v,\ \mathrm{incidence}\ a\ u &\ \Rightarrow\ r(\mathrm{inl}(u,a)) < r(\mathrm{inr}(a,v)).
\end{align}
\end{defn}

Conditions~\eqref{eq:vargt} and~\eqref{eq:facgt} are not chosen for
convenience: they are precisely the membership conditions of
$\mathrm{otherFactors}\ G\ v\ a$ and $\mathrm{otherVars}\ G\ a\ v$, which is
why the next theorem is provable at all.

\begin{thm}[\texttt{bpStep\_rankLocal}]
\label{thm:ranklocal}
For any rank certificate $\mathrm{cert}$ and any factor weight,
$\mathrm{bpStep}\ G\ f$ is rank-local for $\mathrm{cert.rank}$.
\end{thm}

The proof case-splits on the endpoint and unfolds the actual update. At
$\mathrm{inl}(v,a)$ the goal is an equality of products over
$\mathrm{otherFactors}\ G\ v\ a$; a congruence reduces it to an equality of
each factor, and membership in that \texttt{Finset} supplies exactly the
hypothesis of~\eqref{eq:vargt}. At $\mathrm{inr}(a,v)$ the goal is an equality
of sums over assignments, and after congruences on the sum and the inner
product, membership in $\mathrm{otherVars}\ G\ a\ v$ supplies~\eqref{eq:facgt}.
No property of the factor weight is used, so the theorem holds for arbitrary
real weights, including negative ones.

\begin{thm}[\texttt{bp\_converges\_on\_ranked\_acyclic}]
\label{thm:bpconv}
Let $\mathrm{cert}$ be a rank certificate for $G$, let $f$ be any factor weight,
let $m$ be any initial message bundle, and put
$H = \mathrm{convergenceHorizon}\ \mathrm{cert.rank}$. Then every sweep count
$n \geq H$ gives the same bundle as $H$ sweeps, and that bundle satisfies
$\texttt{IsBPFixedPoint}$.
\end{thm}

This follows from Theorem~\ref{thm:converges} and
Theorem~\ref{thm:ranklocal} transported along $\mathrm{flattenMsg}$, whose
injectivity is proven from its left inverse. The quantification is worth
reading carefully: the initial bundle is arbitrary, the factor weight is
arbitrary, the fixed point is an exact equality of bundles, and the bound is
uniform in both.

\subsection{The certificate is not a vacuous hypothesis}

A convergence theorem conditional on a hypothesis is worth only as much as the
hypothesis. Two results pin it down.

\begin{thm}[\texttt{isFourCycleFree\_of\_bpRankCert}]
\label{thm:certc4}
If $G$ admits a rank certificate then $G$ is four-cycle-free.
\end{thm}

A 4-cycle $u\,a\,v\,b$ would give, by two applications each
of~\eqref{eq:vargt} and~\eqref{eq:facgt}, the chain
$r(\mathrm{inr}(a,v)) < r(\mathrm{inl}(v,b)) < r(\mathrm{inr}(b,u)) <
r(\mathrm{inl}(u,a)) < r(\mathrm{inr}(a,v))$, which no rank into $\mathbb{N}$
admits. So a certificate carries genuine loop-freeness content rather than
being satisfiable by fiat.

\begin{prop}[\texttt{bpRankCertStar}, \texttt{bp\_converges\_on\_star}]
\label{prop:star}
The factor graph on two variables and one factor, with the factor incident to
both, carries the explicit certificate $r(\mathrm{inl}(\cdot)) = 0$,
$r(\mathrm{inr}(\cdot)) = 1$. Its bipartite incidence graph is the three-node
path, a genuine tree, and Theorem~\ref{thm:bpconv} applied to it yields a
fixed point for every factor weight and every initial bundle.
\end{prop}

Proposition~\ref{prop:star} is stated as a non-vacuity witness and must be read
with its limitation attached. On that graph~\eqref{eq:facgt} is discharged
substantively, by $0 < 1$ at a real incidence, but~\eqref{eq:vargt} is
discharged vacuously: there is only one factor, so no $b \neq a$ is incident to
either variable. The witness therefore certifies half of
Definition~\ref{def:cert} against a real instance and the other half against an
empty quantifier.

\subsection{Scope}
\label{sec:conv-scope}

Three statements bound what Theorem~\ref{thm:bpconv} establishes, and each is a
gap rather than a caveat.

First, the existence of a rank certificate is a hypothesis and is not derived
from acyclicity. The development defines the bipartite incidence graph, defines
acyclicity and tree-ness through Mathlib's \texttt{SimpleGraph.IsAcyclic} and
\texttt{SimpleGraph.IsTree}, and proves that acyclicity implies
four-cycle-freeness. It does not prove that an acyclic factor graph admits a
certificate. That construction is a finite well-founded one, by leaf stripping
on the forest, and until it lands the theorem reads ``given a topological order
on message endpoints, the iteration freezes'', with the graph theory sitting in
the assumption.

Second, $\mathrm{convergenceHorizon}$ is $\sup(\mathrm{rank}) + 1$ for the
supplied certificate and nothing proves it equals the diameter of the incidence
graph. Since Definition~\ref{def:cert} constrains only the two strict
inequalities, a certificate with unnecessarily spread ranks is still a
certificate and still gives a valid but weaker bound. The diameter reading is
correct for the subtree-depth certificate specifically; that certificate is the
one not yet constructed.

Third, the end-to-end statement that couples convergence to the structural
predicates carries a redundant hypothesis. It takes a certificate and a
tree-simulability hypothesis (tree-ness, non-negative weights, a vanishing
scalar) and returns three conclusions, but all three are reachable from the
certificate together with the two non-structural conjuncts, because
Theorem~\ref{thm:certc4} already supplies four-cycle-freeness from the
certificate alone. The tree hypothesis does no work there.

\figuredeferred{fig:d7-rank-freeze}{message-endpoint ranks on the three-node path with the freeze-out round of each endpoint; blocked on nothing for the star, and on the unbuilt subtree-depth certificate for any deeper tree, so the deeper panel cannot be drawn honestly yet}

\section{A Cram\'er rate for 4-cycle density}
\label{sec:rate}

\section{What the 4-cycle rate does not measure}
\label{sec:limits}

\section{Related work}
\label{sec:related}

\section{What would have to be proved}
\label{sec:outlook}

\appendix

\section{Statements as elaborated}
\label{app:statements}

\section{Axiom closure and reproduction}
\label{app:axioms}

\begin{thebibliography}{99}

\bibitem{TindallMelloFishmanStoudenmireSels2026Science392}
J.~Tindall, A.~F.~Mello, M.~Fishman, E.~M.~Stoudenmire, D.~Sels,
\emph{Dynamics of disordered quantum systems with two- and
three-dimensional tensor networks},
Science \textbf{392}, 868 (2026),
DOI 10.1126/science.adx2728, arXiv:2503.05693.

\bibitem{AntaoSunFumegaLado2026PRL136}
T.~V.~C.~Antão, Y.~Sun, A.~O.~Fumega, J.~L.~Lado,
\emph{Tensor Network Method for Real-Space Topology in Quasicrystal
Chern Mosaics},
Phys.\ Rev.\ Lett.\ \textbf{136}, 156601 (2026),
DOI 10.1103/hhdf-xpwg, arXiv:2506.05230.

\bibitem{YedidiaFreemanWeiss2003}
J.~Yedidia, W.~T.~Freeman, Y.~Weiss,
\emph{Understanding Belief Propagation and Its Generalizations},
in \textit{Exploring Artificial Intelligence in the New Millennium}
(Morgan Kaufmann, 2003), pp.~239--269.

\bibitem{ChertkovChernyak2006LoopSeries}
M.~Chertkov, V.~Y.~Chernyak,
\emph{Loop series for discrete statistical models on graphs},
J.\ Stat.\ Mech.\ (2006) P06009,
DOI 10.1088/1742-5468/2006/06/P06009, arXiv:cond-mat/0603189.

\bibitem{BiancoResta2011}
R.~Bianco, R.~Resta,
\emph{Mapping topological order in coordinate space},
Phys.\ Rev.\ B \textbf{84}, 241106(R) (2011).

\end{thebibliography}

\end{document}
